\documentclass{article}
\usepackage[utf8]{inputenc}
\usepackage{amsmath}
\usepackage{amssymb}
\usepackage{enumitem}
\usepackage{graphicx}

\begin{document}

% Portada
\begin{titlepage}
    \centering
    
\includegraphics[height=2cm]{fing.png}\hfill
    
\includegraphics[height=2cm]{uach.png}
    \\
    \vspace{1.5cm}
    {\Large \textbf{Universidad Autónoma de Chihuahua}}\\[0.2cm]
    {\large Facultad de Ingeniería}\\[0.5cm]
    {\Large \textbf{Métodos Numéricos}}\\[0.5cm]
    {\large Examen Parcialito 2}\\[0.2cm]
    {\large S.S.E.L.S}\\[1cm]

    \vspace{1cm}
    {\large Nombre: \underline{Francisco Javier Ponce Saenz}}\\[0.3cm]
    {\large Matrícula: \underline{a325000}}\\[0.3cm]
    {\large Semestre: \underline{3er Semestre}}\\[0.3cm]
    {\large Fecha: \underline{\hspace{4cm}}}\\
    \vfill
\end{titlepage}

\section*{Examen}

Resuelva el siguiente sistema de ecuaciones lineales simultáneas por los métodos solicitados:

\[
\begin{aligned}
    2x + 3y - 5z + 4t &= 11 \\
    3x - 2y + z  - 2t &= 2 \\
    4x - 6y + 10z - 3t &= 7 \\
    -3x - 3y + 5z - 2t &= -12
\end{aligned}
\]

\begin{enumerate}[label=\Alph*)]
    \item Eliminación de Gauss.\\

    \begin{figure}[h]
        \centering
        \includegraphics[width=0.45\textwidth]{gauss.png}\hfill
        \includegraphics[width=0.45\textwidth]{gauss2.png}
        \\
        \includegraphics[width=0.45\textwidth]{gauss3.png}
        \caption{Pasos de la eliminación de Gauss (parte 1).}
    \end{figure}

    \begin{figure}[h]
        \centering
        \includegraphics[width=0.7\textwidth]{gauss4.png}
        \caption{Paso adicional de la eliminación de Gauss (parte 2).}
    \end{figure}

    \vspace{1cm}

    \item Matriz inversa.\\

    \begin{figure}[h]
        \centering
        \includegraphics[width=0.45\textwidth]{inversa.png}\hfill
        \includegraphics[width=0.45\textwidth]{inversa2.png}
        \\
        \includegraphics[width=0.45\textwidth]{inversa3.png}
        \caption{Pasos para obtener la matriz inversa (parte 1).}
    \end{figure}

    \begin{figure}[h]
        \centering
        \includegraphics[width=0.7\textwidth]{inversa4.png}
        \caption{Paso adicional para obtener la matriz inversa (parte 2).}
    \end{figure}

    \vspace{1cm}
\end{enumerate}

\end{document}
